% Options for packages loaded elsewhere
\PassOptionsToPackage{unicode}{hyperref}
\PassOptionsToPackage{hyphens}{url}
\documentclass[
  ignorenonframetext,
]{beamer}
\newif\ifbibliography
\usepackage{pgfpages}
\setbeamertemplate{caption}[numbered]
\setbeamertemplate{caption label separator}{: }
\setbeamercolor{caption name}{fg=normal text.fg}
\beamertemplatenavigationsymbolsempty
% remove section numbering
\setbeamertemplate{part page}{
  \centering
  \begin{beamercolorbox}[sep=16pt,center]{part title}
    \usebeamerfont{part title}\insertpart\par
  \end{beamercolorbox}
}
\setbeamertemplate{section page}{
  \centering
  \begin{beamercolorbox}[sep=12pt,center]{section title}
    \usebeamerfont{section title}\insertsection\par
  \end{beamercolorbox}
}
\setbeamertemplate{subsection page}{
  \centering
  \begin{beamercolorbox}[sep=8pt,center]{subsection title}
    \usebeamerfont{subsection title}\insertsubsection\par
  \end{beamercolorbox}
}
% Prevent slide breaks in the middle of a paragraph
\widowpenalties 1 10000
\raggedbottom
\AtBeginPart{
  \frame{\partpage}
}
\AtBeginSection{
  \ifbibliography
  \else
    \frame{\sectionpage}
  \fi
}
\AtBeginSubsection{
  \frame{\subsectionpage}
}
\usepackage{iftex}
\ifPDFTeX
  \usepackage[T1]{fontenc}
  \usepackage[utf8]{inputenc}
  \usepackage{textcomp} % provide euro and other symbols
\else % if luatex or xetex
  \usepackage{unicode-math} % this also loads fontspec
  \defaultfontfeatures{Scale=MatchLowercase}
  \defaultfontfeatures[\rmfamily]{Ligatures=TeX,Scale=1}
\fi
\usepackage{lmodern}
\ifPDFTeX\else
  % xetex/luatex font selection
\fi
% Use upquote if available, for straight quotes in verbatim environments
\IfFileExists{upquote.sty}{\usepackage{upquote}}{}
\IfFileExists{microtype.sty}{% use microtype if available
  \usepackage[]{microtype}
  \UseMicrotypeSet[protrusion]{basicmath} % disable protrusion for tt fonts
}{}
\makeatletter
\@ifundefined{KOMAClassName}{% if non-KOMA class
  \IfFileExists{parskip.sty}{%
    \usepackage{parskip}
  }{% else
    \setlength{\parindent}{0pt}
    \setlength{\parskip}{6pt plus 2pt minus 1pt}}
}{% if KOMA class
  \KOMAoptions{parskip=half}}
\makeatother
\setlength{\emergencystretch}{3em} % prevent overfull lines
\providecommand{\tightlist}{%
  \setlength{\itemsep}{0pt}\setlength{\parskip}{0pt}}
\usepackage{bookmark}
\IfFileExists{xurl.sty}{\usepackage{xurl}}{} % add URL line breaks if available
\urlstyle{same}
\hypersetup{
  hidelinks,
  pdfcreator={LaTeX via pandoc}}

\author{\texorpdfstring{}{}}
\date{}

\begin{document}

\begin{frame}[fragile]{Introduction}
\protect\phantomsection\label{introduction}
Rational numbers are dense in the reals, yet the \textbf{quality} of the
best rational approximation at a fixed denominator ceiling varies
sharply across individual \(x\). Number theory classifies extremes:
algebraic irrationals of degree two have bounded partial quotients in
their continued fraction expansion and are \emph{badly approximable} in
the sense of {[}@cassels1957{]}; almost every \(x \in \mathbb{R}\)
satisfies Khinchin-type statistics for the growth of those quotients
{[}@khinchin1964continued{]}. Transcendental constants carry no
finite-degree algebraic constraint, but particular numbers such as
\(\pi\) or \(e\) still admit \textbf{structured} rational approximations
from series, integrals, or independent constructions---so their
finite-\(Q\) behaviour need not resemble a ``generic'' draw from a
simple continuous law.

\begin{block}{Finite \(Q\) as a deliberate lens}
\protect\phantomsection\label{finite-q-as-a-deliberate-lens}
As \(Q\) grows, \(\delta_Q(x)\) tends to \(0\) for every irrational
\(x\), and \(\mu_Q(x)\) admits universal upper bounds tied to
continued-fraction convergents {[}@hurwitz1891{]}. Those limits wash out
differences on a logarithmic scale unless one studies \textbf{rates} or
\textbf{exceptional sets}. By contrast, fixing \(Q\) at a moderate value
(for example \(Q=120\) in the default configuration) asks a different
question: \emph{among all rationals with denominators up to \(Q\), how
small is the best error, and how does that single number compare to the
same functional evaluated on random \(X\)?} The answer is a
\textbf{scalar summary}---an empirical percentile relative to a chosen
reference law---not a Diophantine \emph{type} in the classical sense.
\end{block}

\begin{block}{Statistical comparison protocol}
\protect\phantomsection\label{statistical-comparison-protocol}
The implementation draws a pooled sample \(X_1,\ldots,X_n\) on
\([0,1)\), computes \(\delta_Q(X_i)\) and \(\mu_Q(X_i)\) for each draw,
and records distribution summaries (quantiles, tail percentiles). For
each registered constant \(x^\star\) it stores \(\delta_Q(x^\star)\),
\(\mu_Q(x^\star)\), and the \textbf{empirical rank} (proportion of
reference samples strictly below the constant's value). Optional
\texttt{experiment.q\_sensitivity} repeats the constant table at
additional \(Q\) values \textbf{without resampling} \(X\), so shifts in
rank isolate the effect of enlarging the rational search set.
\end{block}

\begin{block}{Scope of this note}
\protect\phantomsection\label{scope-of-this-note}
We do not prove new Diophantine bounds. The goal is
\textbf{methodological}: connect classical statements about
\(Q\to\infty\) to concrete finite-\(Q\) measurements, document the
numerical chain (three-numerator scan, lattice cross-checks,
convergent-only certificates), and parameterize the study from YAML so
\(Q\), sample sizes, reference components, and \texttt{compare\_mod1}
can change without editing infrastructure code.
\end{block}
\end{frame}

\end{document}
